\documentclass[12pt,a4paper]{article}
% Drake_preamble.tex - LaTeX Preamble Template
% 作者: 謝慕揚, MD, PhD, FESC
% 
% 使用說明:
% 1. 表格請使用 longtable，不要有垂直分隔線 |
% 2. 表格內換行使用 \newline，不要使用 \\
% 3. 藥物名稱、檢驗、檢查請維持英文
% 4. Reference 格式請使用 \href 加上驗證過的 DOI

% Including essential packages for Chinese typesetting
\usepackage{xeCJK}
\usepackage{fontspec}
% Configuring Chinese font with Noto Serif CJK TC, fallback to SimSun
\IfFontExistsTF{Noto Serif CJK TC}{
  \setCJKmainfont{Noto Serif CJK TC}
  \setCJKsansfont{Noto Serif CJK TC}
  \setCJKmonofont{Noto Serif CJK TC}
}{\setCJKmainfont{SimSun}
  \setCJKsansfont{SimSun}
  \setCJKmonofont{SimSun}
}

% Including other necessary packages
\usepackage{geometry}
\usepackage{titlesec}
\usepackage{enumitem}
\usepackage{xcolor}
\usepackage{colortbl}
\usepackage{tcolorbox}
\usepackage{booktabs}
\usepackage{longtable}
\usepackage{array}
\usepackage{graphicx}
\usepackage{tikz}
\usepackage{fancyhdr}
\usepackage{multicol}
\usepackage{datetime2}
\usepackage{hyperref}
\usepackage{mdframed}
\usepackage{amsmath}
\usepackage{amssymb}
\usepackage{multirow}

\hypersetup{
    colorlinks=true,
    linkcolor=emergency_blue,
    citecolor=emergency_blue,
    urlcolor=emergency_blue,
    pdfborder={0 0 0}
}

% Defining page geometry
\geometry{
  top=2.5cm,
  bottom=2.5cm,
  left=2cm,
  right=2cm,
  headheight=25pt
}

% Adjusting topmargin to compensate for increased headheight
\addtolength{\topmargin}{-10pt}

% Defining colors
\definecolor{emergency_red}{RGB}{186,24,27}
\definecolor{emergency_blue}{RGB}{0,114,188}
\definecolor{emergency_gray}{RGB}{120,120,120}
\definecolor{highlight_yellow}{RGB}{255,250,205}

% Setting title formats
\titleformat{\section}[block]{\Large\bfseries\color{emergency_red}}{\thesection}{10pt}{}
\titleformat{\subsection}[block]{\large\bfseries\color{emergency_blue}}{\thesubsection}{8pt}{}
\titleformat{\subsubsection}[block]{\normalsize\bfseries\color{emergency_gray}}{\thesubsubsection}{6pt}{}

% Defining custom environment for important notes
\newmdenv[
    linecolor=emergency_red,
    backgroundcolor=highlight_yellow,
    linewidth=2pt,
    topline=true,
    bottomline=true,
    leftline=true,
    rightline=true,
    innertopmargin=10pt,
    innerbottommargin=10pt,
    innerleftmargin=10pt,
    innerrightmargin=10pt
]{importantbox}

% Defining note command with tcolorbox
\newcommand{\note}[1]{
    \par
    \noindent
    \begin{tcolorbox}[
        colback=emergency_blue!5!white,
        colframe=emergency_blue,
        title=\textbf{重點提醒},
        fonttitle=\bfseries,
        boxrule=0.5pt,
        arc=3pt,
        left=6pt,
        right=6pt,
        top=6pt,
        bottom=6pt
    ]
    #1
    \end{tcolorbox}
}

% Key findings box
\newcommand{\keyfinding}[1]{
    \par
    \noindent
    \begin{tcolorbox}[
        colback=yellow!10!white,
        colframe=orange!75!black,
        title=\textbf{核心發現摘要},
        fonttitle=\bfseries,
        boxrule=0.5pt,
        arc=3pt
    ]
    #1
    \end{tcolorbox}
}

% Defining custom date format for ISO 8601
\DTMnewdatestyle{iso}{
  \renewcommand{\DTMdisplaydate}[4]{\number##1-\number##2-\number##3}
  \renewcommand{\DTMDisplaydate}[4]{\number##1-\number##2-\number##3}
}
\DTMsetdatestyle{iso}
\newcommand{\mydate}{\DTMtoday}


% Custom header for this document
\pagestyle{fancy}
\fancyhf{}
\fancyhead[L]{\textcolor{emergency_red}{PCRonline 教學講義}}
\fancyhead[R]{\textcolor{emergency_blue}{Tricuspid Intervention 2025}}
\fancyfoot[C]{\textcolor{emergency_gray}{\thepage}}
\renewcommand{\headrulewidth}{1pt}
\renewcommand{\headrule}{\hbox to\headwidth{\color{emergency_red}\leaders\hrule height \headrulewidth\hfill}}

\begin{document}

%============================================================
% Title Page
%============================================================
\begin{center}
{\LARGE\bfseries\textcolor{emergency_red}{PCRonline}}\\[0.5cm]
{\Large\textbf{The Top Ten Impactful Papers from 2025 on Tricuspid Intervention We Almost Missed}}\\[0.3cm]
{\large 2025年三尖瓣介入治療十大重要文獻回顧}\\[0.5cm]
{\normalsize \href{https://www.pcronline.com/News/Whats-new-on-PCRonline/2026/The-top-ten-impactful-papers-on-tricuspid-intervention-we-almost-missed}{PCRonline. 13 Jan 2026. Review by Alessandro Sticchi}}\\[0.3cm]
{\normalsize 整理：謝慕揚 MD, PhD, FESC \quad 日期：\mydate}
\end{center}

\vspace{0.5cm}

%============================================================
\section{前言}
%============================================================

經導管三尖瓣介入治療 (Transcatheter tricuspid intervention) 近年快速發展，已從可行性驗證進入臨床應用階段。除了 TRILUMINATE、TRISCEND 等里程碑試驗外，2025年有許多較少被關注但極具價值的研究，探討了介入後的生理機制、併發症、病人表型及非預期後果。

\keyfinding{
\textbf{五大核心發現：}
\begin{enumerate}[nosep]
    \item \textbf{傳導障礙}：TTVR 後新發傳導障礙達 44\%，14-19\% 需永久性心律調節器 (PPM)
    \item \textbf{導線交互作用}：T-TEER 對 RV lead 相對安全；TTVR 後 31\% 出現導線異常
    \item \textbf{亞臨床血栓}：CT 偵測 HALT 達 27\%，嚴重瓣葉活動度下降 20\%
    \item \textbf{雙心室重塑}：RVEDV 下降 65\%、LVEDV 上升 17\%、effective RVEF 改善 65\%
    \item \textbf{病人表型}：高劑量利尿劑組死亡率 44\% vs 11\%（HR 2.50）
\end{enumerate}
}

%============================================================
\section{十大重要研究摘要}
%============================================================

\begin{longtable}{p{0.5cm}p{3.5cm}p{3cm}p{3.5cm}p{4cm}}
\toprule
\textbf{\#} & \textbf{研究 (第一作者)} & \textbf{族群與設計} & \textbf{主要數據} & \textbf{核心訊息} \\
\midrule
\endfirsthead
\toprule
\textbf{\#} & \textbf{研究 (第一作者)} & \textbf{族群與設計} & \textbf{主要數據} & \textbf{核心訊息} \\
\midrule
\endhead
\bottomrule
\endfoot

1 & Le Ruz et al. (2025) & 70 TTVR 病人\newline 單中心研究 & NOCD 44.3\%\newline PPM 14\%\newline RV strain: NOCD+ 29.7\% vs 25.1\% & TTVR 後傳導障礙常見，與 RV 功能及 device-septal interaction 相關 \\
\midrule

2 & Storozhenko et al. (2026)\newline \textbf{Tri-LEAD Study} & 146 T-TEER 病人\newline 含 RV leads & Lead dysfunction 6.8\%\newline 無 fracture 或再介入\newline 追蹤約 2 年 & T-TEER 時 RV lead entrapment 中期追蹤相對安全 \\
\midrule

3 & Abbasi et al. (2025) & 32 TTVR 病人\newline 含 CIEDs & Lead abnormalities 31\%\newline Insulation breach 13\%\newline Revision 3\%\newline >80\% 在 90 天內 & TTVR 後 lead jailing 常致早期電學變化，需密切監測 \\
\midrule

4 & Peigh et al. (2025) & 52 TTVR 病人\newline Entrapped RV leads & Major lead events 7.7\%\newline 早期 fracture/\newline dislodgement & TTVR 後導線併發症不常見但不可忽視，尤其早期 \\
\midrule

5 & Kempton et al. (2025) & 45 TTVR 病人\newline 常規 CT 追蹤 & HALT 27\%\newline Severe RLM 20\%\newline 較高 gradient\newline NYHA 改善較差 & TTVR 後亞臨床瓣膜血栓常見，TTE 難以偵測 \\
\midrule

6 & Le Ruz et al. (2025)\newline \textbf{Remodeling} & 80 TTVR 病人\newline Echo + CT & RVEDV $-$65\%\newline LVEDV +17\%\newline Effective RVEF +65\% & TTVR 引發快速且顯著的雙心室逆向重塑 \\
\midrule

7 & Marx et al. (2026) & 225 TTVr 病人 & HDD 死亡率 44\% vs 11\%\newline HFH 22\% vs 10\%\newline HR 2.50 & 高利尿劑劑量代表進展性 TR 表型，預後較差 \\
\midrule

8 & Chen et al. (2025) & 159 TTVR 病人\newline 大 vs 小 annulus & Procedural success\newline 81\% vs 97\%\newline Clinical success\newline 74\% vs 95\% & 大 annulus 降低手術成功率，但不影響早期安全性 \\
\midrule

9 & Lurz et al. (2025)\newline \textbf{Trillium Device} & 20 heterotopic TTVR 病人 & Device success 100\%\newline CVP $-$3 mmHg\newline NYHA 改善 & 異位 cross-caval TTVR 為無法接受其他治療的病人提供症狀緩解 \\
\midrule

10 & Angellotti et al. (2025)\newline \textbf{Real-world EVOQUE} & 176 real-world\newline EVOQUE TTVR & TR $\leq$ mild 98.4\%\newline NYHA I-II: 20\%$\rightarrow$80\%\newline PPM 18.9\% & 真實世界 TTVR 高度有效，但 pacemaker 負擔顯著 \\

\end{longtable}

%============================================================
\section{關鍵主題深入分析}
%============================================================

\subsection{傳導障礙與心律調節器需求}

\begin{itemize}[leftmargin=*]
    \item TTVR 後新發傳導障礙 (NOCD) 發生率高達 \textbf{44.3\%}，其中 \textbf{14-19\%} 需植入 PPM
    \item 機制相關因素：device 與 membranous septum 接觸、RV strain 較高者風險增加
    \item 基線已有傳導疾病者風險更高
    \item \textbf{臨床建議}：術前應建立「傳導風險評估框架」，而非將 PPM 視為可接受的附帶效應
\end{itemize}

\subsection{導線交互作用 (Lead-Device Interaction)}

\note{
\textbf{T-TEER vs TTVR 的差異：}
\begin{itemize}[nosep]
    \item T-TEER：Lead dysfunction 僅 6.8\%，無 fracture 或需再介入，中期安全
    \item TTVR：Lead abnormalities 31\%，insulation breach 13\%，需 revision 3\%
    \item 多數異常出現在術後 60-90 天內
\end{itemize}
}

\begin{itemize}[leftmargin=*]
    \item 「Lead jailing」並非單一策略，而是需要個別化評估的過程
    \item 建議：
    \begin{enumerate}[nosep]
        \item 術前標準化電生理評估
        \item 明確告知病人相關風險
        \item 術後前 60-90 天加強監測
    \end{enumerate}
\end{itemize}

\subsection{亞臨床瓣膜血栓 (Subclinical Valve Thrombosis)}

\begin{itemize}[leftmargin=*]
    \item CT 可偵測到經胸超音波 (TTE) 無法發現的瓣膜異常：
    \begin{itemize}[nosep]
        \item HALT (Hypoattenuated Leaflet Thickening): 27\%
        \item Severe RLM (Reduced Leaflet Motion): 20\%
    \end{itemize}
    \item HALT/RLM 與較高的 transvalvular gradient 及較差的症狀改善相關
    \item \textbf{臨床意涵}：症狀改善不佳、gradient 上升或血流動力學異常的病人，應考慮 CT 追蹤
\end{itemize}

\subsection{雙心室逆向重塑 (Biventricular Reverse Remodeling)}

\begin{itemize}[leftmargin=*]
    \item TTVR 術後 RV 容積快速下降，LV 充盈增加：
    \begin{itemize}[nosep]
        \item RVEDV 下降 65\%
        \item LVEDV 上升 17\%
        \item Effective RVEF 改善 65\%
        \item RV-PA coupling 改善 20\%
    \end{itemize}
    \item 這重新定義了 TR 治療的「成功」標準：從單純的瓣膜導向目標轉變為雙心室及全身性目標
\end{itemize}

\subsection{病人表型與預後 (Patient Phenotype)}

\begin{importantbox}
\textbf{高利尿劑劑量 (High Diuretic Dose, HDD) 的預後意義}

\begin{itemize}[nosep]
    \item 定義：Furosemide-equivalent $>$ 95 mg/day
    \item 1 年死亡率：HDD 44\% vs 非 HDD 11\%
    \item 心衰住院：HDD 22\% vs 非 HDD 10\%
    \item HR 2.50
\end{itemize}

即使 HDD 組仍有症狀改善且無需利尿劑加量，但預後明顯較差。\\
\textbf{利尿劑劑量升高可能是介入時機的警訊，也可能定義「為時已晚」的界線。}
\end{importantbox}

\subsection{解剖學考量：大 Annulus}

\begin{itemize}[leftmargin=*]
    \item Large annulus 與 small annulus 比較：
    \begin{itemize}[nosep]
        \item Procedural success: 81\% vs 97\%
        \item Clinical success: 74\% vs 95\%
    \end{itemize}
    \item 大 annulus 降低手術成功率，但不影響早期安全性
    \item 需發展針對大 annulus 的技術改良或替代策略
\end{itemize}

\subsection{異位置換策略：Heterotopic TTVR}

\begin{itemize}[leftmargin=*]
    \item Trillium device：Cross-caval 異位置換
    \item 治療目標不同：處理 TR 的全身靜脈後果，而非瓣膜本身
    \item 結果：
    \begin{itemize}[nosep]
        \item Device success 100\%
        \item CVP 下降 3 mmHg
        \item NYHA 改善
    \end{itemize}
    \item \textbf{臨床定位}：為無法接受 orthotopic TTVR 的病人提供替代選擇
\end{itemize}

%============================================================
\section{未來方向與臨床整合}
%============================================================

\begin{center}
\begin{tcolorbox}[
    colback=emergency_blue!5!white,
    colframe=emergency_blue,
    title=\textbf{三尖瓣介入治療的典範轉移},
    fonttitle=\bfseries,
    width=0.95\textwidth
]
\begin{enumerate}
    \item \textbf{從手術技巧到疾病管理}：整合解剖、電生理、影像監測、充血生理及技術選擇

    \item \textbf{術前評估標準化}：
    \begin{itemize}[nosep]
        \item 傳導風險評估
        \item CIED 導線評估與電生理會診
        \item 利尿劑劑量作為疾病嚴重度指標
        \item Annulus 大小對 device 選擇的影響
    \end{itemize}

    \item \textbf{術後監測強化}：
    \begin{itemize}[nosep]
        \item 前 60-90 天密切監測導線功能
        \item 症狀改善不佳者考慮 CT 評估 HALT/RLM
    \end{itemize}

    \item \textbf{治療策略分層}：
    \begin{itemize}[nosep]
        \item Orthotopic TTVR：適合能耐受且可受益於完全 TR 消除的病人
        \item Heterotopic TTVR：優先目標為全身充血減壓及症狀緩解的病人
    \end{itemize}
\end{enumerate}
\end{tcolorbox}
\end{center}

%============================================================
\section{縮寫對照表}
%============================================================

\begin{longtable}{p{3.5cm}p{11cm}}
\toprule
\textbf{縮寫} & \textbf{全名} \\
\midrule
\endfirsthead
\midrule
\textbf{縮寫} & \textbf{全名} \\
\midrule
\endhead
\bottomrule
\endfoot

TTVR & Transcatheter Tricuspid Valve Replacement（經導管三尖瓣置換術）\\
TTVr & Transcatheter Tricuspid Valve Repair（經導管三尖瓣修補術）\\
T-TEER & Tricuspid Transcatheter Edge-to-Edge Repair（三尖瓣經導管緣對緣修補術）\\
TR & Tricuspid Regurgitation（三尖瓣逆流）\\
RV & Right Ventricle（右心室）\\
RA & Right Atrium（右心房）\\
NOCD & New-Onset Conduction Disturbances（新發傳導障礙）\\
PPM & Permanent Pacemaker Implantation（永久性心律調節器植入）\\
CIED & Cardiac Implantable Electronic Device（心臟植入式電子裝置）\\
HALT & Hypoattenuated Leaflet Thickening（瓣葉低衰減增厚）\\
RLM & Reduced Leaflet Motion（瓣葉活動度下降）\\
CT & Computed Tomography（電腦斷層）\\
HFH & Heart Failure Hospitalization（心衰竭住院）\\
HDD & High Diuretic Dose（高劑量利尿劑）\\
CVP & Central Venous Pressure（中心靜脈壓）\\
RVEDV & Right Ventricular End-Diastolic Volume（右心室舒張末期容積）\\
LVEDV & Left Ventricular End-Diastolic Volume（左心室舒張末期容積）\\
RVEF & Right Ventricular Ejection Fraction（右心室射出分率）\\

\end{longtable}

%============================================================
\section{參考文獻}
%============================================================

\begin{enumerate}
    \item Le Ruz R, Nazif T, George I, et al. New-Onset Conductance Disturbances After Transcatheter Tricuspid Valve Replacement: A Mechanistic Assessment. \href{https://doi.org/10.1016/j.jcin.2025.07.039}{\textit{JACC Cardiovasc Interv}. 2025;18:2569-79.}

    \item Storozhenko T, Russo G, Vanderheyden M, et al. Tricuspid Right Ventricular Lead Entrapment in Transcatheter Tricuspid Interventions: The Tri-LEAD Study. \href{https://doi.org/10.1016/j.jacep.2025.11.003}{\textit{JACC Clin Electrophysiol}. 2025.}

    \item Abbasi M, Killu AM, Van Niekerk C, et al. Device-lead abnormalities and function after transcatheter tricuspid valve replacement. \href{https://doi.org/10.1093/europace/euaf219}{\textit{Europace}. 2025;27:euaf219.}

    \item Peigh G, Al-Kazaz M, Davidson LJ, et al. Outcomes of Entrapped Right Ventricular Pacing or Defibrillator Leads Following Transcatheter Tricuspid Valve Replacement. \href{https://doi.org/10.1016/j.jcin.2025.05.042}{\textit{JACC Cardiovasc Interv}. 2025;18:1762-72.}

    \item Kempton H, Stolz L, Stocker T, et al. Hypoattenuated Leaflet Thickening and Reduced Leaflet Motion After Transcatheter Tricuspid Valve Replacement. \href{https://doi.org/10.1016/j.jcin.2025.10.026}{\textit{JACC Cardiovasc Interv}. 2025.}

    \item Le Ruz R, Agarwal V, George I, et al. Cardiac Remodeling After Transcatheter Tricuspid Valve Replacement. \href{https://doi.org/10.1016/j.jcin.2025.10.023}{\textit{JACC Cardiovasc Interv}. 2025.}

    \item Marx L, Finke K, Althoff J, et al. Relationship between Diuretic Dose and Outcome Following Transcatheter Tricuspid Valve Repair for Severe Tricuspid Regurgitation. \href{https://doi.org/10.1016/j.cjca.2025.12.014}{\textit{Can J Cardiol}. 2025.}

    \item Chen F, Yang X, Qiao F, et al. Comparison of Early Outcomes Following Transjugular Transcatheter Tricuspid Valve Replacement Between Large vs Small Annulus. \href{https://doi.org/10.1016/j.jcin.2025.11.013}{\textit{JACC Cardiovasc Interv}. 2025.}

    \item Lurz P, Kresoja KP, Besler C, et al. Heterotopic Crosscaval Transcatheter Tricuspid Valve Replacement for Patients With Tricuspid Regurgitation: The Trillium Device. \href{https://doi.org/10.1016/j.jcin.2025.04.036}{\textit{JACC Cardiovasc Interv}. 2025;18:1425-34.}

    \item Angellotti D, Mattig I, Samim D, et al. Early Outcomes of Real-World Transcatheter Tricuspid Valve Replacement. \href{https://doi.org/10.1016/j.jcin.2025.06.002}{\textit{JACC Cardiovasc Interv}. 2025;18:1896-909.}
\end{enumerate}

\end{document}
